\section{System Overview}

\subsection{Overall Workflow}
The proposed workflow links heterogeneous traffic inputs to downstream ETA-oriented evaluation through a single artifact ecosystem. At a high level, the system begins from one of two data-construction paths: a simulation-derived path built on controlled SUMO scenarios, and a real-trajectory path built on imported GPS-like traces aligned to a road network. Although these paths differ substantially in observability, preprocessing, and error sources, both terminate in a common graph-based representation designed for route-aware ETA tasks. This shared representation can then be exercised through a web-based evaluation environment in which a compatible predictor is attached, trip-level outcomes are registered, and statistical analyses are produced.

\subsection{Desktop Dataset Construction Environment}
Dataset construction is operationalized through the \DesktopTool{}, an interactive, user-facing desktop environment rather than a collection of disconnected preprocessing scripts. On the simulation side, the desktop application supports scenario definition, control of simulation behavior, visual validation of semantic elements such as zones, landmarks, and vehicle classes, and export of route-aware graph snapshots with adjustable temporal granularity. It also allows users to tune aspects of scenario realism, including configurable noise such as random untracked trips. On the trajectory side, the same environment supports map download, network preparation, trajectory repair, route alignment, and batch export into the common representation.

\subsection{Shared ETA-Oriented Representation}
The conceptual center of the system is the common ETA-oriented representation produced by both construction paths. In this formulation, roads are modeled as edges, while junctions and vehicles are modeled as nodes. The representation is explicitly route-aware: each vehicle is associated with a known or inferred route, so that the resulting data more closely resembles a navigation-style setting than a generic traffic snapshot. It is also dynamic in both attributes and structure, as traffic state evolves over time and graph relations change with vehicle movement. Although this representation is primarily designed for vehicle-level ETA tasks, it is sufficiently rich to support other downstream views such as road-level load estimation or area-level traffic analysis.

\subsection{Web-Based Testing Environment}
The \WebApp{} provides the downstream evaluation layer of the workflow. Rather than stopping at offline dataset generation, the system supports interactive testing sessions in which a compatible ETA predictor is attached to the produced representation and exercised under controlled conditions. Prediction outcomes are registered in a database, associated with trip-level information, and exposed through statistical summaries, plots, and report generation. In this way, the web environment serves as the bridge between representation design and measurable downstream use.

\subsection{The Predictive Model in the Overall System}
The predictive model occupies a deliberately bounded role in the architecture. It is included as a representation-compatible validation component rather than as the main contribution of the paper. The framework is designed so that compatible predictors can be attached, replaced, and tested within the same evaluation environment, including by external users of the open-source release. The cloud-deployed instance discussed later is therefore best understood as one practical demonstrator configuration of a broader dataset-construction and testing workflow.
